% Class-compatibility preamble retained for a reader excerpt of
%   Dave Witte Morris, "A Cohomological Proof that Real Representations of
%   Semisimple Lie Algebras Have Q-Forms", SIGMA 11 (2015) 034,
%   DOI 10.3842/SIGMA.2015.034, CC BY 4.0.
%
% The source pins the publisher class sigma, which is not part of the delivery
% baseline, so this file supplies the amsart baseline and then copies the author's
% own declaration block mechanically and unchanged:
%   - pinned-original-source.tex lines 11-31: the author's DeclareMathOperator
%     list, the equation numbering and the author's newtheorem declarations;
%   - pinned-original-source.tex lines 35-43, which the author writes immediately
%     after \begin{document} in the publisher layout: the author's own \Res,
%     \ResOne and \Lie macros and \allowdisplaybreaks.
% The publisher header commands of the sigma class (\FirstPageHeading,
% \ArticleName and the like) and the macro \PaperNumber they require are not
% reproduced, because this excerpt carries neither the journal front matter nor
% anything else that uses them.
% No mathematical macro, statement, symbol or assumption is introduced, changed
% or removed by this file.  The author's full DeclareMathOperator list of pinned
% source lines 11-19 is carried, including GL, SL, SO and SU: these four are not
% supplied by the delivery baseline, so leaving them out makes the author's own
% notation (for instance GL_n(Q) in the corollary on lattices) undefined.
\documentclass[10pt,a4paper]{amsart}
\usepackage{amssymb,amsmath,amsthm}
\usepackage{enumerate}
\usepackage{hyperref}

\DeclareMathOperator{\Cliff}{Clif\/f}
\DeclareMathOperator{\Gal}{Gal}
\DeclareMathOperator{\GL}{GL}
\DeclareMathOperator{\Inn}{Inn}
\DeclareMathOperator{\Mat}{Mat}
\DeclareMathOperator{\Out}{Out}
\DeclareMathOperator{\SL}{SL}
\DeclareMathOperator{\SO}{SO}
\DeclareMathOperator{\SU}{\mathbf{SU}}
\numberwithin{equation}{section}

\newtheorem{prop}{Proposition}[section]
\newtheorem{cor}[prop]{Corollary}
\newtheorem{lem}[prop]{Lemma}
{\theoremstyle{definition}
\newtheorem{defn}[prop]{Definition}
\newtheorem{notation}[prop]{Notation}
\newtheorem{rem}[prop]{Remark}
\newtheorem{ques}[prop]{Question}
\newtheorem{eg}[prop]{Example}
\newtheorem{erratum}[prop]{Erratum}}

\newcommand{\Res}[1]{\mathop{\mathbf{R}_{#1}}}
\newcommand{\ResOne}[1]{\mathop{\mathbf{R}^{(1)}_{#1}}}
\newcommand{\Lie}[1]{\mathfrak{\lowercase{#1}}}
\allowdisplaybreaks
\usepackage{cleveref}
\crefname{originaldefinition}{original definition}{original definitions}
\Crefname{originaldefinition}{Original definition}{Original definitions}
\makeatletter
\newlabel{BetaDefn}{{1.1}{source}{}{sourcepointer-1}{}}
\newlabel{BetaDefn@cref}{{[originaldefinition][0][]1.1}{[1][0][]source}{}{}{}}
\makeatother
\crefname{originallemma}{original lemma}{original lemmas}
\Crefname{originallemma}{Original lemma}{Original lemmas}
\makeatletter
\newlabel{RunivIff}{{2.1}{source}{}{sourcepointer-2}{}}
\newlabel{RunivIff@cref}{{[originallemma][0][]2.1}{[1][0][]source}{}{}{}}
\makeatother
\crefname{originaldefinition}{original definition}{original definitions}
\Crefname{originaldefinition}{Original definition}{Original definitions}
\makeatletter
\newlabel{RunivIff-D}{{4}{source}{}{sourcepointer-3}{}}
\newlabel{RunivIff-D@cref}{{[originaldefinition][0][]4}{[1][0][]source}{}{}{}}
\makeatother
\crefname{originallemma}{original lemma}{original lemmas}
\Crefname{originallemma}{Original lemma}{Original lemmas}
\makeatletter
\newlabel{SameFixed}{{2.3}{source}{}{sourcepointer-4}{}}
\newlabel{SameFixed@cref}{{[originallemma][0][]2.3}{[1][0][]source}{}{}{}}
\makeatother
\title{A simple non-absolutely-simple algebra is R-universal exactly when it is the restriction of scalars of a strongly inner algebra over an imaginary quadratic extension}
\author{Dave Witte Morris}
\date{doi:10.3842/SIGMA.2015.034; CC BY 4.0}
\begin{document}
\maketitle
\noindent\textit{Editorial note.} All delivered mathematical text is the author's own, copied from the pinned publisher source sources/algebra-fl02-a-cohomological-proof-that-real-re/pinned-original-source.tex; no statement or proof step was rewritten, reordered or removed.\par\medskip
\noindent\textit{Editorial note.} Preamble provenance: the source pins the publisher class sigma, which is not part of the delivery baseline, so this card uses sources/algebra-fl02-a-cohomological-proof-that-real-re/retained-preamble.tex, which supplies the amsart baseline and then carries the author's own declarations from the pinned source (its header records the exact line ranges). The author's full DeclareMathOperator list is carried, including GL, SL, SO and SU, because the delivery baseline does not provide those four operator names. No macro was invented and no mathematical symbol was changed.\par\medskip
\noindent\textit{Editorial note.} Bibliography provenance: the citations of the counted proof are resolved with sources/algebra-fl02-a-cohomological-proof-that-real-re/bibliography.tex, an excerpt of the author's own thebibliography environment (pinned source lines 954-1016) copied character for character with all twelve entries kept; no entry was rewritten, reordered, merged, split or invented and no citation key was changed.\par\medskip
\noindent\textit{Editorial note.} External references are navigation pointers only; every deliberately kept step of the author's argument stays inside the counted spans.\par\medskip
\noindent\textit{Editorial note.} Selection provenance: this card was opened under the file-based control contract (CONTROL.md), with the mechanical pre-card gate run on its statement and proof interval before the card was written.\par\medskip
\noindent\textit{Allowed standard prerequisites (editorial).}\begin{itemize}
\item The description of a simple non-absolutely-simple algebra as the restriction of scalars of an absolutely simple algebra over a finite extension. Used at the start of the counted proof (source lines 857-858), attributed by the author to section 3.1.2 of the cited paper Tits-Classification. Supplies the shape of the algebra that the criterion is stated in terms of.
\item The relation between the beta invariants of an algebra and those of its restriction of scalars. Used inside the counted proof where the author compares the invariant of the restricted weight with the invariant of the original one (source lines 859-877). Transfers the vanishing condition for the beta invariant across the restriction of scalars.
\end{itemize}\medskip
\ifdefined\hypertarget\hypertarget{sourcepointer-1}{}\fi
\noindent\textit{Source reference.} Original definition 1.1: sources/algebra-fl02-a-cohomological-proof-that-real-re/pinned-original-source.tex, source lines 135--135. Navigation pointer only: the author's definition of the beta invariant, used by the counted proof in its notation.\par\medskip
\ifdefined\hypertarget\hypertarget{sourcepointer-2}{}\fi
\noindent\textit{Source reference.} Original lemma 2.1: sources/algebra-fl02-a-cohomological-proof-that-real-re/pinned-original-source.tex, source lines 416--416. Navigation pointer only: the counted proof invokes the author's equivalence for R-universality; that lemma has its own card in this lane and its content is declared as a named standard prerequisite.\par\medskip
\ifdefined\hypertarget\hypertarget{sourcepointer-3}{}\fi
\noindent\textit{Source reference.} Original definition 4: sources/algebra-fl02-a-cohomological-proof-that-real-re/pinned-original-source.tex, source lines 432--432. Navigation pointer only: the labelled item of that equivalence used by the counted proof.\par\medskip
\ifdefined\hypertarget\hypertarget{sourcepointer-4}{}\fi
\noindent\textit{Source reference.} Original lemma 2.3: sources/algebra-fl02-a-cohomological-proof-that-real-re/pinned-original-source.tex, source lines 471--471. Navigation pointer only: the author's lemma on the minimal inner extension; it has its own card in this lane and its content is declared as a named standard prerequisite.\par\medskip
\begin{prop}
Let $\Lie G$ be a~simple Lie algebra over~$\mathbb{Q}$ that is not absolutely simple.
Then $\Lie G$ is $\mathbb{R}$-universal if and only if $\Lie G = \Res{L/\mathbb{Q}} \Lie G'$ is the restriction of
scalars of a~strongly inner, absolutely simple Lie algebra~$\Lie G'$ over an imaginary quadratic extension~$L$
of~$\mathbb{Q}$.
\end{prop}

\begin{proof}
($\Rightarrow$) We have $\Lie G = \Res{L/\mathbb{Q}} \Lie G'$, for some absolutely simple Lie algebra~$\Lie G'$ over
some f\/inite extension~$L$ of~$\mathbb{Q}$~\cite[\S~3.1.2]{Tits-Classification}.
From Lemma~\ref{SameFixed}, we see that if $\Lie G$ is $\mathbb{R}$-universal, then the extension~$L$ must be imaginary
quadratic.

Let $\lambda'$ be a~nonzero dominant weight of~$\Lie G'$, and let $\lambda = \Res{L/\mathbb{Q}} \lambda'$ be the
corresponding weight of~$\Lie G$.
Then~$\lambda$ is not f\/ixed by the $*$-action, so we have $D_{\Lie g, F}(\lambda) = D_{\Lie g, L}(\lambda) = D_{\Lie g',
L}(\lambda)$.
(The f\/inal equality is because $\Lie G = \Res{L/\mathbb{Q}} \Lie G'$.) On the other hand, since all Lie algebras are
split over~$\mathbb{C}$, we have $D_{\Lie g, \mathbb{R}}(\lambda) = D_{\Lie g, \mathbb{C}}(\lambda) = \mathbb{C}$.
Hence, we see from Lemma~\ref{RunivIff}(\ref{RunivIff-D}) that if $\Lie G$ is $\mathbb{R}$-universal, then $D_{\Lie g', L}(\lambda)
= L$.
Since this is true for every dominant weight~$\lambda$ of~$\Lie G'$, we conclude that $\Lie G'$ is strongly inner.

($\Leftarrow$) Note that, since $\Lie G'$ is strongly inner, the simply connected group~$\mathbf{G}'$ corresponding
to~$\Lie G'$ may be written as $\mathbf{G}' = {}^{\zeta} \mathbf{G}^s$, where $\mathbf{G}^s$ is split (and simply
connected) and $\zeta \in H^1(L; \mathbf{G}^s)$.
By Restriction of Scalars, then $\mathbf{G} = {}^{\xi} \mathbf{G}^q$, where $\mathbf{G}^q = \Res{L/\mathbb{Q}}
\mathbf{G}^s$ is quasi-split and $\xi \in H^1(\mathbb{Q}; \mathbf{G}^q)$.
Then, in the notation of Def\/inition~\ref{BetaDefn}, we have $\delta_*[\xi] = 0$, so $\beta_{\Lie G, F}(\lambda) = 0$ for every
$*$-invariant weight~$\lambda$.
It is therefore easy to see that $\Lie G$ is $\mathbb{R}$-universal.
\end{proof}

%% Bibliography excerpt retained for reader excerpts of
%   Dave Witte Morris, "A Cohomological Proof that Real Representations of Semisimple Lie Algebras Have Q-Forms",
%   SIGMA 11 (2015) 034, DOI 10.3842/SIGMA.2015.034.
%%
%% Every entry below is copied character for character from the author's own
% thebibliography environment of the pinned publisher source (lines 954-1016).
%% No entry was rewritten, reordered, merged, split or invented, and no citation
% key was changed.  All 12 entries are kept so that any citation of a delivered
%% excerpt resolves without guessing an entry.
\begin{thebibliography}{99}

\footnotesize\itemsep=0pt

\bibitem{DemazureGrothendieck-SGA3}
Demazure M., Grothendieck A., Sch\'emas en groups.~III. Structure des sch\'emas
  en groupes reductifs, S\'eminaire de g\'eom\'etrie alg\'ebrique du {B}ois
  {M}arie, 1962/1964, available at
  \url{http://library.msri.org/books/sga/sga/pdf/sga3-3.pdf}.

\bibitem{Harder-BerichtNeuere}
Harder G., Bericht \"uber neuere {R}esultate der {G}aloiskohomologie
  halbeinfacher {G}ruppen, \textit{Jber. Deutsch. Math.-Verein.} \textbf{70}
  (1968), 182--216, available at
  \url{http://www.digizeitschriften.de/en/dms/img/?PPN=PPN37721857X_0070&DMDID=dmdlog17}.

\bibitem{Kneser-LectGalCohoClassical}
Kneser M., Lectures on {G}alois cohomology of classical groups, \textit{Tata
  Institute of Fundamental Research Lectures on Mathematics}, Vol.~47, Tata
  Institute of Fundamental Research, Bombay, 1969, available at
  \url{http://www.math.tifr.res.in/~publ/ln/tifr47.pdf}.

\bibitem{KnusEtAl-BookOfInvols}
Knus M.A., Merkurjev A., Rost M., Tignol J.P., The book of involutions,
  \textit{American Mathematical Society Colloquium Publications}, Vol.~44,
  Amer. Math. Soc., Providence, RI, 1998.

\bibitem{Lam-IntroQuadForms}
Lam T.Y., Introduction to quadratic forms over f\/ields, \textit{Graduate Studies
  in Mathematics}, Vol.~67, Amer. Math. Soc., Providence, RI, 2005.

\bibitem{Morris-RRepsQForms}
Morris D., Real representations of semisimple {L}ie algebras have {${\mathbb
  Q}$}-forms, in Algebraic Groups and Arithmetic, Tata Institute of Fundamental
  Research, Mumbai, 2004, 469--490.

\bibitem{PlatonovRapinchukBook}
Platonov V., Rapinchuk A., Algebraic groups and number theory, \textit{Pure and
  Applied Mathematics}, Vol.~139, Academic Press, Inc., Boston, MA, 1994.

\bibitem{PrasadRapinchuk-Prescribed}
Prasad G., Rapinchuk A., On the existence of isotropic forms of semi-simple
  algebraic groups over number f\/ields with prescribed local behavior,
  \href{http://dx.doi.org/10.1016/j.aim.2006.01.001}{\textit{Adv. Math.}} \textbf{207} (2006), 646--660.

\bibitem{Raghunathan-ArithLattsSSGrps}
Raghunathan M.S., Arithmetic lattices in semisimple groups, \href{http://dx.doi.org/10.1007/BF02967980}{\textit{Proc.
  Indian Acad. Sci. Math. Sci.}} \textbf{91} (1982), 133--138.

\bibitem{Tits-Classification}
Tits J., Classif\/ication of algebraic semisimple groups, in Algebraic {G}roups
  and {D}iscontinuous {S}ubgroups ({P}roc. {S}ympos. {P}ure {M}ath., {B}oulder,
  {C}olo., 1965), Amer. Math. Soc., Providence, R.I., 1966, 33--62.

\bibitem{Tits-RepLinIrred}
Tits J., Repr\'esentations lin\'eaires irr\'eductibles d'un groupe r\'eductif
  sur un corps quelconque, \href{http://dx.doi.org/10.1515/crll.1971.247.196}{\textit{J.~Reine Angew. Math.}} \textbf{247} (1971),
  196--220.

\bibitem{Tits-StronglyInner}
Tits J., Strongly inner anisotropic forms of simple algebraic groups,
  \href{http://dx.doi.org/10.1016/0021-8693(90)90201-X}{\textit{J.~Algebra}} \textbf{131} (1990), 648--677.


\end{thebibliography}
\end{document}
